\begin{proof}[Proof of \cref{obj:lem:normalized-affine-testing}]
Set
\[
c=\frac{10^{-20}}{128}>0,
\]
and fix public parameters satisfying the hypotheses. Throughout the proof, the quantities
\[
S,\ N,\ \ell,\ x,\ u,\ w,\ M,\ L,\ B,\ \alpha,\ b_0,\ K_*
\]
have exactly the values in the affine branch of \cref{obj:def:prior-handle}.

\begin{enumerate}
\item
We first establish support and the common marginal for an arbitrary admissible certificate \(\sigma\). The tuning formulas give
\[
128N\le M\le256N,\qquad
8\ell\le L\le9\ell,\qquad
B,\alpha,b_0>0,\qquad K_*b_0\le\frac1{100}.
\]
Indeed, \(M=128(n+N)\) and \(n\le N\). Also \(\ell>1\), so
\[
8\ell\le\lceil8\ell\rceil<8\ell+1\le9\ell.
\]
The last bound follows by multiplying
\[
K_*\le\left\lfloor\frac1{100b_0}\right\rfloor
\le\frac1{100b_0}
\]
by \(b_0\). In particular, \(L\ge2\), and \(K_*\le d-1\).

Write the certificate and its variation measure as
\[
\sigma=\sum_{i=1}^{L+2}\omega_i\delta_{v_i},
\qquad
\eta=|\sigma|=\sum_{i=1}^{L+2}|\omega_i|\delta_{v_i}.
\]
For configurations involving a zero-weight node, extend the polar sign by
\[
h(v_i)=
\begin{cases}
1,&\omega_i\ge0,\\
-1,&\omega_i<0.
\end{cases}
\]
Thus \(|\omega_i|h(v_i)=\omega_i\), including when \(\omega_i=0\).

The one-coordinate latent space and its probability law are
\[
\Lambda=\{(0,0)\}\cup\{(v_i,1):1\le i\le L+2\},
\]
\[
\begin{aligned}
\nu_{\mathrm{cell}}\{(v_i,1)\}
&=\frac{\alpha}{S_1(v_i)}|\omega_i|,\\
\nu_{\mathrm{cell}}\{(0,0)\}
&=1-\sum_{i=1}^{L+2}\frac{\alpha}{S_1(v_i)}|\omega_i|.
\end{aligned}
\]
The branch flag distinguishes the filler state \((0,0)\) from a core state at the node zero. Since \(S_1(v)\ge\alpha>0\),
\[
0\le\sum_i\frac{\alpha}{S_1(v_i)}|\omega_i|
\le\sum_i|\omega_i|=1,
\]
so these masses are nonnegative and sum to one.

For a coordinate \((V,c_i)\) with this law,
\[
\mathbb EV
=\int v\frac{\alpha}{S_1(v)}\,d\eta(v)
\le\frac{\alpha}{1-\epsilon},
\]
because \(v\alpha/S_1(v)\le\alpha/(1-\epsilon)\). Consequently,
\[
\bar w=b_0+\mathbb EV
\le b_0+\frac{\alpha}{1-\epsilon}
\le\frac43b_0,
\]
where \(\alpha=\epsilon b_0\) and \(\epsilon\le1/4\). Hence, for every latent configuration,
\[
p_*=1-K_*\bar w
\ge1-\frac1{75}>\frac12,
\qquad
Q=p_*+\sum_{j=1}^{K_*}(b_0+V_j)>\frac12.
\]
These assertions also cover an empty rare-cell set, in which case \(p_*=Q=1\).

For every \(v\in[0,B]\), the raw arm masses satisfy
\[
\begin{aligned}
S_1(v)+S_0(v)&=b_0+v=w_*(v),\\
S_1(v)-\epsilon w_*(v)&=(1-2\epsilon)v\ge0,\\
(1-\epsilon)w_*(v)-S_1(v)&=(1-2\epsilon)b_0\ge0.
\end{aligned}
\]
Both raw arms are positive. The assigned conditional means belong to \([0,1]\), and summing the raw complete-record table over all atoms gives
\[
p_*+\sum_{j=1}^{K_*}w_*(V_j)=Q.
\]
Division by \(Q\) therefore produces probability tables. Their occupied rare cells satisfy the overlap inequalities above; the reservoir has propensity \(1/2\); and the remaining cells have mass zero. Thus
\[
P_\iota(\boldsymbol\lambda)\in\mathcal M_{d,\epsilon},
\qquad
\iota\in\{0,1\},\quad
\boldsymbol\lambda=((V_j,c_j))_{j=1}^{K_*}\in\Lambda^{K_*},
\]
by \cref{obj:def:model}.

Summing out the outcome coordinate gives, under either hypothesis,
\[
(P_\iota(\boldsymbol\lambda))_{XA}(j,a)=
\begin{cases}
S_a(V_j)/Q,&1\le j\le K_*,\\
p_*/(2Q),&j=K_*+1,\\
0,&j>K_*+1.
\end{cases}
\]
In particular,
\[
(P_1(\boldsymbol\lambda))_{XA}
=(P_0(\boldsymbol\lambda))_{XA}
\quad\text{for every }\boldsymbol\lambda.
\]
Define the common product latent law and its pushforward priors by
\[
\nu_{\mathrm{lat}}=\nu_{\mathrm{cell}}^{\otimes K_*},
\qquad
\Pi_\iota^\sigma=\operatorname{Law}_{\nu_{\mathrm{lat}}}
\bigl(P_\iota(\boldsymbol\lambda)\bigr).
\]
Pushing the pointwise equality through this common latent law yields
\[
\operatorname{Law}_{\Pi_1^\sigma}(P_{XA})
=\operatorname{Law}_{\Pi_0^\sigma}(P_{XA}).
\]
% lean: CausalSmith.Stat.AnnotationRarearmFrontier.affine_prior_support_and_common_marginal

\item
The assumed inequalities
\[
S\ge\exp(4096),\qquad x^2>S^{-1}
\]
select the affine branch of \cref{obj:def:prior-handle}. The preceding equality applies to its selected certificate, and hence
\[
\operatorname{Law}_{\Pi_1}(P_{XA})
=\operatorname{Law}_{\Pi_0}(P_{XA}).
\]
For the risk bound, fix an admissible certificate. Such a certificate exists because \(L\ge2\), \(B>0\), and \(0<\epsilon\le1/4\): \cref{obj:lem:shrinking-cone-dual} supplies the strictly ordered nodes, variation mass one, vanishing moments through degree \(L\), and reciprocal gap
\[
\int\frac{\alpha}{\alpha+(1-\epsilon)v}\,d\sigma(v)\ge\frac18.
\]
% lean: CausalSmith.Stat.AnnotationRarearmFrontier.affine_selectedPrior_common_marginal
% lean: CausalSmith.Stat.AnnotationRarearmFrontier.affine_schedule_dual

\item
Define the raw targets, their means, their gap, and the normalized targets by
\[
\begin{aligned}
Z_\iota^{\mathrm{tar}}(\boldsymbol\lambda)
&=\sum_{j=1}^{K_*}w_*(V_j)\mu_1^{(\iota)}(V_j,c_j),\\
t_\iota^{\mathrm{tar}}
&=\int Z_\iota^{\mathrm{tar}}\,d\nu_{\mathrm{lat}},\\
\Delta_{\mathrm{tar}}&=t_1^{\mathrm{tar}}-t_0^{\mathrm{tar}},\\
\tau_\iota(\boldsymbol\lambda)
&=\tau(P_\iota(\boldsymbol\lambda))
=\frac{Z_\iota^{\mathrm{tar}}(\boldsymbol\lambda)}{Q(\boldsymbol\lambda)}.
\end{aligned}
\]
The last identity follows from \cref{obj:def:target}: the control means and the reservoir outcome means are zero. Pointwise,
\[
0\le Z_\iota^{\mathrm{tar}}
\le\sum_{j=1}^{K_*}w_*(V_j)\le Q,
\qquad
0\le\tau_\iota\le1.
\]

The difference of the one-cell target means is
\[
\mathbb E_{\nu_{\mathrm{cell}}}
\left[w_*(V)\{\mu_1^{(1)}(V,c_i)-\mu_1^{(0)}(V,c_i)\}\right]
=\int w_*(v)\frac{\alpha}{S_1(v)}\,d\sigma(v).
\]
Here the filler contributes zero, while multiplying the core probability by the polar sign restores the signed weight. The identity
\[
w_*(v)\frac{\alpha}{S_1(v)}
=\frac{\alpha}{1-\epsilon}
+\frac{b_0(1-2\epsilon)}{1-\epsilon}
  \frac{\alpha}{S_1(v)}
\]
and the zero-th moment of \(\sigma\) give
\[
\int w_*(v)\frac{\alpha}{S_1(v)}\,d\sigma(v)
=\frac{b_0(1-2\epsilon)}{1-\epsilon}
 \int\frac{\alpha}{S_1(v)}\,d\sigma(v)
\ge\frac{b_0}{12},
\]
since \((1-2\epsilon)/(1-\epsilon)\ge2/3\). Additivity of expectation now gives
\[
\Delta_{\mathrm{tar}}\ge\frac{K_*b_0}{12}.
\]

We next calibrate this gap in terms of \(x\). The tuning formulas imply
\[
b_0=\frac1{100000M\epsilon L}
\ge\frac1{230400000N\epsilon\ell}.
\]
As \(S\ge1\), we have \(M\epsilon\ge128S\) and \(L\ge8\), so \(b_0\le1/100\). Thus
\[
\left\lfloor\frac1{100b_0}\right\rfloor\ge1.
\]
The defining floor inequality then gives
\[
\frac1{100b_0}
<
\left\lfloor\frac1{100b_0}\right\rfloor+1
\le2\left\lfloor\frac1{100b_0}\right\rfloor,
\qquad
\left\lfloor\frac1{100b_0}\right\rfloor b_0\ge\frac1{200}.
\]
Together with \(d-1\ge d/2\), this yields
\[
(d-1)b_0\ge\frac{x}{460800000}.
\]
If the alphabet cap determines \(K_*\), this is the required mass bound. If the floor cap determines \(K_*\), the bound \(1/200\) applies. Therefore
\[
K_*b_0\ge\frac{\min\{x,1\}}{460800000},
\qquad
\Delta_{\mathrm{tar}}\ge10^{-10}\min\{x,1\}>0.
\]
In particular, \(K_*\ge1\) in the present regime.
% lean: CausalSmith.Stat.AnnotationRarearmFrontier.affine_raw_target_mean_gap_uniform

\item
Consider the following raw Poisson experiment. Conditional on
\(\boldsymbol\lambda\) and hypothesis \(\iota\), the complete-record atom counts and auxiliary atom counts are mutually independent, with respective means
\[
uQ\,P_\iota(\boldsymbol\lambda)(j,a,y),
\qquad
wQ\,(P_\iota(\boldsymbol\lambda))_{XA}(j,a).
\]
We use the Poisson convention
\[
\operatorname{Poi}(\lambda_{\mathrm P})\{k\}
=e^{-\lambda_{\mathrm P}}\frac{\lambda_{\mathrm P}^k}{k!},
\qquad \lambda_{\mathrm P}\ge0,\quad k\in\mathbb N,
\]
with \(0^0=1\).

In one rare cell, combine the two control-channel counts. The resulting four-count mixture is
\[
\begin{aligned}
\nu_\iota^{\mathrm{cell}}
=\int_{\Lambda}
&\operatorname{Poi}\bigl(uS_1(v)\mu_1^{(\iota)}(v,c_i)\bigr)
\otimes
\operatorname{Poi}\bigl(uS_1(v)(1-\mu_1^{(\iota)}(v,c_i))\bigr)\\
&\otimes\operatorname{Poi}\bigl(wS_1(v)\bigr)
\otimes\operatorname{Poi}\bigl(MS_0(v)\bigr)
\,d\nu_{\mathrm{cell}}(v,c_i).
\end{aligned}
\]
The first two coordinates count treated complete outcomes equal to one and zero, respectively.

When both complete treated counts are zero, the two mixture probabilities agree. When both are positive, both mixture probabilities vanish. For
\(\rho,\gamma,\kappa\in\mathbb N\), define the signed discrepancy
\[
D_{\rho,\gamma,\kappa}^{\mathrm{cnt}}
=
\int e^{-M(b_0+v)}
\frac{(uS_1(v))^{\rho+1}(wS_1(v))^\gamma(MS_0(v))^\kappa}
     {(\rho+1)!\gamma!\kappa!}
\frac{\alpha}{S_1(v)}\,d\sigma(v).
\]
Expanding the conditional Poisson probabilities and the core mixture gives
\[
\begin{aligned}
\nu_1^{\mathrm{cell}}\{(\rho+1,0,\gamma,\kappa)\}
-\nu_0^{\mathrm{cell}}\{(\rho+1,0,\gamma,\kappa)\}
&=D_{\rho,\gamma,\kappa}^{\mathrm{cnt}},\\
\nu_1^{\mathrm{cell}}\{(0,\rho+1,\gamma,\kappa)\}
-\nu_0^{\mathrm{cell}}\{(0,\rho+1,\gamma,\kappa)\}
&=-D_{\rho,\gamma,\kappa}^{\mathrm{cnt}}.
\end{aligned}
\]
Indeed, the filler has treated mean zero under both hypotheses. On the core branch, the difference between the two deterministic outcome assignments is the polar sign, and
\(|\omega_i|h(v_i)=\omega_i\). These are all the possible nonzero discrepancies.
% lean: CausalSmith.Stat.AnnotationRarearmFrontier.affineCellPoissonPredictive_equation_six

\item
We bound the sum of these discrepancies by moment cancellation. Define
\[
\begin{aligned}
\varphi_{\rho,\gamma,\kappa}(v)
&=S_1(v)^{\rho+\gamma}S_0(v)^\kappa,\\
\beta_{\rho,\gamma,\kappa}^{\mathrm{cnt}}
&=e^{-Mb_0}
  \frac{u^\rho w^\gamma M^\kappa}{(\rho+1)!\gamma!\kappa!},\\
a_j^{\mathrm{cnt}}
&=\sum_{\rho,\gamma,\kappa\ge0}
  \beta_{\rho,\gamma,\kappa}^{\mathrm{cnt}}
  [v^j]\varphi_{\rho,\gamma,\kappa}(v),
\qquad j\in\mathbb N.
\end{aligned}
\]
Here \([v^j]\varphi(v)=\varphi^{(j)}(0)/j!\) denotes the coefficient of \(v^j\). Factoring one treated intensity gives
\[
D_{\rho,\gamma,\kappa}^{\mathrm{cnt}}
=u\alpha\,\beta_{\rho,\gamma,\kappa}^{\mathrm{cnt}}
\int e^{-Mv}\varphi_{\rho,\gamma,\kappa}(v)\,d\sigma(v).
\]
Every coefficient of \(\varphi_{\rho,\gamma,\kappa}\) is nonnegative.

For \(j\le L\), the matched moments imply
\[
\int e^{-Mv}v^j\,d\sigma(v)
=
\int\left(e^{-Mv}
-\sum_{k=0}^{L-j}\frac{(-Mv)^k}{k!}\right)v^j\,d\sigma(v).
\]
Bounding the exponential remainder on \([0,B]\), and using
\(\|\sigma\|_{\mathrm{TV}}=1\), gives
\[
\left|\int e^{-Mv}v^j\,d\sigma(v)\right|
\le B^j\sum_{\substack{k\ge0\\k+j>L}}\frac{(MB)^k}{k!}.
\]
For \(j>L\), the full absolute exponential series gives the same inequality, with every \(k\ge0\) included. Expanding the polynomial therefore yields
\[
\left|\int e^{-Mv}\varphi_{\rho,\gamma,\kappa}(v)\,d\sigma(v)\right|
\le
\sum_{j\ge0}[v^j]\varphi_{\rho,\gamma,\kappa}(v)\,
B^j\sum_{\substack{k\ge0\\k+j>L}}\frac{(MB)^k}{k!}.
\]

To sum over the counts, use \(1/(\rho+1)!\le1/\rho!\). Coefficientwise, the nonnegative count series obey
\[
\begin{aligned}
\sum_{\rho\ge0}\frac{u^\rho S_1(v)^\rho}{(\rho+1)!}
&\preceq e^{uS_1(v)},\\
\sum_{\gamma\ge0}\frac{w^\gamma S_1(v)^\gamma}{\gamma!}
&=e^{wS_1(v)},\\
\sum_{\kappa\ge0}\frac{M^\kappa S_0(v)^\kappa}{\kappa!}
&=e^{MS_0(v)},
\end{aligned}
\]
where \(\preceq\) means that each coefficient on the left is at most the corresponding coefficient on the right. Thus
\[
0\le a_j^{\mathrm{cnt}}
\le e^{-Mb_0}[v^j]e^{M(S_1(v)+S_0(v))}
=[v^j]e^{Mv}
=\frac{M^j}{j!}.
\]
This calculation cancels the common baseline using
\(S_1(v)+S_0(v)=b_0+v\).
The dominating series is finite:
\[
\sum_{j\ge0}\frac{M^jB^j}{j!}
\sum_{k\ge0}\frac{(MB)^k}{k!}
=e^{2MB}.
\]
Consequently the nonnegative coefficient sums can be interchanged, and the signed discrepancies are absolutely summable. We obtain
\[
\begin{aligned}
\sum_{\rho,\gamma,\kappa\ge0}
|D_{\rho,\gamma,\kappa}^{\mathrm{cnt}}|
&\le
u\alpha\sum_{j\ge0}a_j^{\mathrm{cnt}}B^j
\sum_{\substack{k\ge0\\k+j>L}}\frac{(MB)^k}{k!}\\
&\le
u\alpha\sum_{\substack{j,k\ge0\\j+k>L}}
\frac{(MB)^{j+k}}{j!k!}\\
&=
u\alpha\sum_{k_{\mathrm{tail}}>L}\frac{(2MB)^{k_{\mathrm{tail}}}}{k_{\mathrm{tail}}!}.
\end{aligned}
\]
The last equality uses
\[
\sum_{j=0}^{k_{\mathrm{tail}}}\frac1{j!(k_{\mathrm{tail}}-j)!}=\frac{2^{k_{\mathrm{tail}}}}{k_{\mathrm{tail}}!}.
\]
% lean: CausalSmith.Stat.AnnotationRarearmFrontier.affine_signed_raw_count_tail_bound

\item
For probability measures, use the total variation convention
\[
\operatorname{TV}(\nu,\nu')
=\sup_{\mathscr E\ \mathrm{measurable}}
|\nu(\mathscr E)-\nu'(\mathscr E)|.
\]
On a countable space this equals one half of the sum of absolute singleton discrepancies. The two disjoint patterns in the preceding calculation therefore give
\[
\operatorname{TV}(\nu_1^{\mathrm{cell}},\nu_0^{\mathrm{cell}})
\le u\alpha\sum_{k_{\mathrm{tail}}>L}\frac{(2MB)^{k_{\mathrm{tail}}}}{k_{\mathrm{tail}}!}.
\]

The tuning gives \(2MB=L/500\). For an integer \(k_{\mathrm{tail}}>L\),
\[
\frac{(L/500)^{k_{\mathrm{tail}}}}{k_{\mathrm{tail}}!}
\le\frac{(k_{\mathrm{tail}}/500)^{k_{\mathrm{tail}}}}{k_{\mathrm{tail}}!}
\le\left(\frac{e}{500}\right)^{k_{\mathrm{tail}}}
\le e^{-5k_{\mathrm{tail}}}.
\]
The middle inequality follows by retaining the \(k_{\mathrm{tail}}\)-th term of the positive series for \(e^{k_{\mathrm{tail}}}\). For the last inequality, \(e\le11/4\) gives
\(e^6\le(11/4)^6<500\). Since \(e^{-5}\le1/2\),
\[
\sum_{k_{\mathrm{tail}}>L}\frac{(L/500)^{k_{\mathrm{tail}}}}{k_{\mathrm{tail}}!}
\le\frac{e^{-5(L+1)}}{1-e^{-5}}
\le2e^{-5L}.
\]
Moreover,
\[
K_*u\alpha=128S(K_*b_0)\le\frac{32}{25}S,
\]
and
\[
Se^{-5L}
\le e^{\ell-5L}
\le e^{-39\ell}
\le e^{-8}\le\frac1{256}.
\]
Here \(S\le e^\ell\), \(L\ge8\ell\), \(\ell\ge1\), and
\(e^8\ge2^8=256\). Hence
\[
K_*\operatorname{TV}(\nu_1^{\mathrm{cell}},\nu_0^{\mathrm{cell}})
\le K_*u\alpha\sum_{k_{\mathrm{tail}}>L}\frac{(2MB)^{k_{\mathrm{tail}}}}{k_{\mathrm{tail}}!}
\le\frac1{100}<\frac1{64}.
\]

We now restore the full observation channels. Given a combined control count \(\kappa\), its complete-channel count has conditional distribution
\[
\Pr(\kappa_{\mathrm L}=k\mid\kappa)
=\binom{\kappa}{k}\left(\frac uM\right)^k
                    \left(\frac wM\right)^{\kappa-k},
\qquad 0\le k\le\kappa.
\]
This follows directly from
\[
\begin{aligned}
&\operatorname{Poi}(uS_0(v))\{k\}
 \operatorname{Poi}(wS_0(v))\{\kappa-k\}\\
&\hspace{1cm}=
\operatorname{Poi}(MS_0(v))\{\kappa\}
\binom{\kappa}{k}\left(\frac uM\right)^k
                    \left(\frac wM\right)^{\kappa-k}.
\end{aligned}
\]
The splitting law is independent of the latent coordinate and hypothesis. Applying it cannot increase total variation. Independence of the latent coordinates makes the rare-cell mixtures a product, whose distance is bounded by the sum of the cell distances; this follows by telescoping the difference of the product mass functions. The reservoir counts form a common independent factor, and the remaining cells have zero counts. Consequently the full two-channel histogram mixtures have distance below \(1/64\).

Uniformly ordering each channel's histogram is another common probability kernel. Its conditional law is the ordered Poisson experiment from the normalized table. To see this, for a probability mass function \(p_{\mathrm{mark}}\) on a finite alphabet \(\mathscr Z\), put
\[
\kappa_{\mathrm{tot}}=\sum_{z\in\mathscr Z}\kappa_z.
\]
The independent-count probabilities factor as
\[
\prod_{z\in\mathscr Z}
e^{-\lambda_{\mathrm P}p_{\mathrm{mark}}(z)}
\frac{(\lambda_{\mathrm P}p_{\mathrm{mark}}(z))^{\kappa_z}}{\kappa_z!}
=
e^{-\lambda_{\mathrm P}}
\frac{\lambda_{\mathrm P}^{\kappa_{\mathrm{tot}}}}{\kappa_{\mathrm{tot}}!}
\frac{\kappa_{\mathrm{tot}}!}{\prod_z\kappa_z!}
\prod_zp_{\mathrm{mark}}(z)^{\kappa_z}.
\]
Thus, conditional on the latent table, the ordered channels have the following law:
\[
\begin{aligned}
J_{\mathrm L}&\sim\operatorname{Poi}(uQ),&
\mathcal L_{\mathrm{raw}}\mid J_{\mathrm L}=k
&\sim P_\iota(\boldsymbol\lambda)^{\otimes k},\\
J_{\mathrm A}&\sim\operatorname{Poi}(wQ),&
\mathcal X_{\mathrm{raw}}\mid J_{\mathrm A}=k
&\sim(P_\iota(\boldsymbol\lambda))_{XA}^{\otimes k}.
\end{aligned}
\]
Both channels are independent conditional on the latent table. Their domains are
\[
\mathcal L_{\mathrm{raw}}\in
\bigsqcup_{k\ge0}
\bigl(\{1,\ldots,d\}\times\{0,1\}^2\bigr)^k,
\qquad
\mathcal X_{\mathrm{raw}}\in
\bigsqcup_{k\ge0}
\bigl(\{1,\ldots,d\}\times\{0,1\}\bigr)^k,
\]
with the empty tuple at length zero.

Append the independent seed \(U\sim\mathrm{Uniform}[0,1]\), and define
\[
\mathcal Q_\iota^{\mathrm{raw}}
=\operatorname{Law}(\mathcal L_{\mathrm{raw}},
                    \mathcal X_{\mathrm{raw}},U\mid\iota),
\]
where the latent table is mixed according to \(\nu_{\mathrm{lat}}\). Common kernels and a common independent seed preserve the distance bound, so
\[
\operatorname{TV}(\mathcal Q_1^{\mathrm{raw}},
                  \mathcal Q_0^{\mathrm{raw}})<\frac1{64}.
\]
% lean: CausalSmith.Stat.AnnotationRarearmFrontier.affineOrderedRawKernel_randomized_tv_small

\item
We next bound the normalized target's deviation from its raw mean. Since \(0\le V\le B\),
\[
\mathbb EV^2\le B\mathbb EV
\le\frac{B\alpha}{1-\epsilon}
\le\frac43B\alpha.
\]
For either hypothesis,
\[
0\le w_*(V)\mu_1^{(\iota)}(V,c_i)\le w_*(V),
\qquad
\mathbb Ew_*(V)^2
\le2b_0^2+2\mathbb EV^2
\le2b_0^2+\frac83B\alpha.
\]
Independence and the inequality
\(\operatorname{Var}(Z)\le\mathbb EZ^2\) consequently give
\[
\operatorname{Var}(Z_\iota^{\mathrm{tar}})
\le K_*\left(2b_0^2+\frac83B\alpha\right),
\qquad
\mathbb EQ=1,
\qquad
\operatorname{Var}(Q)\le\frac43K_*B\alpha.
\]
For the last two identities, use \(p_*=1-K_*\bar w\) and
\(\operatorname{Var}(b_0+V)=\operatorname{Var}(V)\).

The normalization identity is
\[
\tau_\iota-t_\iota^{\mathrm{tar}}
=(Z_\iota^{\mathrm{tar}}-t_\iota^{\mathrm{tar}})
+\tau_\iota(1-Q).
\]
Since \(0\le\tau_\iota\le1\), squaring and integrating yields
\[
\begin{aligned}
\mathbb E_{\nu_{\mathrm{lat}}}
(\tau_\iota-t_\iota^{\mathrm{tar}})^2
&\le2\operatorname{Var}(Z_\iota^{\mathrm{tar}})
   +2\operatorname{Var}(Q)\\
&\le4K_*b_0^2+8K_*B\alpha.
\end{aligned}
\]
Markov's inequality at the positive radius \(\Delta_{\mathrm{tar}}/8\), together with
\(\Delta_{\mathrm{tar}}\ge K_*b_0/12\), gives
\[
\begin{aligned}
\nu_{\mathrm{lat}}\left\{
|\tau_\iota-t_\iota^{\mathrm{tar}}|
\ge\frac{\Delta_{\mathrm{tar}}}{8}\right\}
&\le
\frac{64(4K_*b_0^2+8K_*B\alpha)}
     {\Delta_{\mathrm{tar}}^2}\\
&\le\frac{36864+7372800\epsilon^2L^2}{K_*}\\
&\le\frac{2^{23}(1+\epsilon^2L^2)}{K_*}.
\end{aligned}
\]
The middle estimate uses the tuning identity
\[
B\alpha=100\epsilon^2L^2b_0^2.
\]
% lean: CausalSmith.Stat.AnnotationRarearmFrontier.affine_normalized_target_gap_tail

\item
To calibrate this tail bound, define the ordinary logarithmic scale
\[
\ell_{\mathrm{plain}}=\log S.
\]
The information-scale assumption implies
\[
\ell_{\mathrm{plain}}\ge4096,\qquad
\ell\le\ell_{\mathrm{plain}}+1,\qquad
L\le10\ell_{\mathrm{plain}}.
\]
For the second inequality, \(S\ge2\) gives \(e+S\le eS\); the third follows from \(L\le9\ell\).

The floor-mass bound already established implies
\[
\left\lfloor\frac1{100b_0}\right\rfloor
\ge\frac1{200b_0}=500M\epsilon L
\ge64000SL.
\]
Also \(x^2>S^{-1}\) and \(x>0\) give
\[
d>\frac{N\epsilon\ell}{\sqrt S}\ge\sqrt S\,\ell.
\]
Thus the two caps defining \(K_*\) satisfy
\[
d-1\ge\frac d2>\frac{\sqrt S\,\ell}{2}
\ge\frac{\sqrt S\,L}{18},
\qquad
\left\lfloor\frac1{100b_0}\right\rfloor
\ge64000SL\ge\frac{\sqrt S\,L}{18},
\]
and hence
\[
K_*\ge\frac{\sqrt S\,L}{18}.
\]

The required exponential comparison follows from
\[
2^{50}\le e^{\ell_{\mathrm{plain}}/4},
\qquad
\frac{\ell_{\mathrm{plain}}}{4}
\le e^{\ell_{\mathrm{plain}}/4}.
\]
The first uses \(e>2\) and
\(50\le\ell_{\mathrm{plain}}/4\); the second follows from
\(1+t\le e^t\). Multiplying gives
\[
e^{\ell_{\mathrm{plain}}/2}
\ge2^{48}\ell_{\mathrm{plain}}
\ge180\cdot2^{30}\ell_{\mathrm{plain}}.
\]
Since \(\sqrt S=e^{\ell_{\mathrm{plain}}/2}\),
\[
K_*\ge10\cdot2^{30}\ell_{\mathrm{plain}}L.
\]
Finally,
\[
1\le\ell_{\mathrm{plain}}L,\qquad
\epsilon^2L^2\le\frac{L^2}{16}
\le\frac{10}{16}\ell_{\mathrm{plain}}L,
\]
so
\[
1+\epsilon^2L^2\le10\ell_{\mathrm{plain}}L,
\qquad
K_*\ge2^{30}(1+\epsilon^2L^2).
\]
Substituting into the preceding tail bound proves, for both hypotheses,
\[
\nu_{\mathrm{lat}}\left\{
|\tau_\iota-t_\iota^{\mathrm{tar}}|
\ge\frac{\Delta_{\mathrm{tar}}}{8}\right\}
\le\frac1{128}.
\]
% lean: CausalSmith.Stat.AnnotationRarearmFrontier.affine_normalized_target_tail_small

\item
Let \(\widehat\tau_{\mathrm{raw}}\) be any Borel rule on the ordered raw channels and seed. Define its two Bayes squared risks by
\[
\mathcal A_\iota(\widehat\tau_{\mathrm{raw}})
=
\int
\mathbb E\left[
\bigl(\widehat\tau_{\mathrm{raw}}
(\mathcal L_{\mathrm{raw}},\mathcal X_{\mathrm{raw}},U)
-\tau_\iota(\boldsymbol\lambda)\bigr)^2
\mid\iota,\boldsymbol\lambda
\right]d\nu_{\mathrm{lat}}(\boldsymbol\lambda).
\]
Threshold the rule at
\[
t_{\mathrm{mid}}^{\mathrm{tar}}
=\frac{t_0^{\mathrm{tar}}+t_1^{\mathrm{tar}}}{2},
\]
and define the test errors
\[
\begin{aligned}
e_0^{\mathrm{test}}
&=\mathcal Q_0^{\mathrm{raw}}
  \{\widehat\tau_{\mathrm{raw}}\ge t_{\mathrm{mid}}^{\mathrm{tar}}\},\\
e_1^{\mathrm{test}}
&=\mathcal Q_1^{\mathrm{raw}}
  \{\widehat\tau_{\mathrm{raw}}<t_{\mathrm{mid}}^{\mathrm{tar}}\}.
\end{aligned}
\]
The definition of total variation gives
\[
e_0^{\mathrm{test}}+e_1^{\mathrm{test}}
\ge1-\operatorname{TV}(\mathcal Q_0^{\mathrm{raw}},
                       \mathcal Q_1^{\mathrm{raw}})
\ge1-\frac1{64}.
\]
On a test error with
\(|\tau_\iota-t_\iota^{\mathrm{tar}}|\le\Delta_{\mathrm{tar}}/8\),
the estimation error is at least \(3\Delta_{\mathrm{tar}}/8\). The two target-tail bounds therefore imply
\[
\mathcal A_0(\widehat\tau_{\mathrm{raw}})
+\mathcal A_1(\widehat\tau_{\mathrm{raw}})
\ge
\left(\frac{3\Delta_{\mathrm{tar}}}{8}\right)^2
\left(e_0^{\mathrm{test}}+e_1^{\mathrm{test}}-\frac2{128}\right).
\]
Consequently,
\[
\max_{\iota\in\{0,1\}}\mathcal A_\iota(\widehat\tau_{\mathrm{raw}})
\ge
\frac12\left(\frac{3\Delta_{\mathrm{tar}}}{8}\right)^2
\left(1-\frac1{64}-\frac2{128}\right)
=\frac{279}{4096}\Delta_{\mathrm{tar}}^2
\ge\frac{\Delta_{\mathrm{tar}}^2}{64}.
\]
If either Bayes risk is infinite, this conclusion holds directly. The common seed is already part of the experiment, so this argument includes randomized rules.
% lean: CausalSmith.Stat.AnnotationRarearmFrontier.affineRawTaggedKernel_bayes_lower

\item
We transfer this bound to an arbitrary clipped Borel rule \(T\) in the original experiment. First we control prefix failures. If
\(J\sim\operatorname{Poi}(\lambda_{\mathrm P})\) and
\(\lambda_{\mathrm P}\ge64k\), exponential Markov gives
\[
\Pr(J<k)
\le\Pr(J\le k)
\le4^k\mathbb E4^{-J}
=\exp\left(k\log4-\frac34\lambda_{\mathrm P}\right)
\le e^{-31k},
\]
using \(\log4\le3\). The event \(J<0\) is empty.

Since \(Q>1/2\), the conditional means satisfy
\[
uQ\ge64n,\qquad wQ\ge64N.
\]
For the retention event
\[
\mathcal E_{\mathrm{ret}}=\{J_{\mathrm L}\ge n,\ J_{\mathrm A}\ge m\},
\]
the complete-channel tail is bounded at \(n\), while
\(\{J_{\mathrm A}<m\}\subseteq\{J_{\mathrm A}<N\}\). Hence, uniformly in the latent configuration and hypothesis,
\[
\Pr(\mathcal E_{\mathrm{ret}}^c\mid\iota,\boldsymbol\lambda)
\le e^{-31n}+e^{-31N}\le2e^{-31n}.
\]
When \(m=0\), the auxiliary failure event is empty.

We calibrate this last quantity against the target gap. Since
\(\epsilon\le1/4\), \(n\ge4S\). For \(S\ge1\),
\[
S\le e^{124(S-1)},
\qquad
Se^{-124S}\le e^{-124}.
\]
Also \(e^{124}\ge2^{124}\), so
\[
2e^{-31n}\le2e^{-124S}\le\frac{10^{-24}}{S}.
\]
If \(x\le1\), the affine-regime inequality gives
\((\min\{x,1\})^2=x^2>S^{-1}\). If \(x>1\), then
\((\min\{x,1\})^2=1\ge S^{-1}\). The calibrated mean gap consequently satisfies
\[
\Delta_{\mathrm{tar}}^2
\ge10^{-20}(\min\{x,1\})^2
\ge\frac{10^{-20}}S.
\]
It follows that
\[
\Pr(\mathcal E_{\mathrm{ret}}^c\mid\iota,\boldsymbol\lambda)
\le2e^{-31n}
\le\frac{\Delta_{\mathrm{tar}}^2}{128}.
\]

Define the original rule's risk and its worst-case risk by
\[
\begin{aligned}
\mathcal R_P(T)
&=\mathbb E_{P^{\otimes n}\otimes P_{XA}^{\otimes m}
                  \otimes\mathrm{Uniform}[0,1]}
  [(T(\mathcal D,U)-\tau(P))^2],\\
\mathcal R_{\max}(T)
&=\sup_{P\in\mathcal M_{d,\epsilon}}\mathcal R_P(T).
\end{aligned}
\]
Clipping and \(\tau(P)\in[-1,1]\) give
\[
0\le\mathcal R_P(T)\le4.
\]
Thus this supremum is finite, and every constructed table's risk is at most \(\mathcal R_{\max}(T)\).

The transferred rule is
\[
\widetilde T(\mathcal L_{\mathrm{raw}},\mathcal X_{\mathrm{raw}},U)
=
\begin{cases}
T(\mathcal L_{\mathrm{raw}}[1:n],
  \mathcal X_{\mathrm{raw}}[1:m],U),
&\mathcal E_{\mathrm{ret}},\\
0,&\mathcal E_{\mathrm{ret}}^c.
\end{cases}
\]
The brackets denote the first indicated number of ordered records, with the empty prefix when that number is zero. This is a Borel rule.

Conditional on the latent table and the channel totals, the records are independent draws from the normalized laws. Thus, on retention, their prefixes have exactly the original two-channel product law, independently of the seed. In particular,
\[
\begin{aligned}
&\mathbb E\left[
(\widetilde T-\tau_\iota)^2\mathbf1_{\mathcal E_{\mathrm{ret}}}
\mid\iota,\boldsymbol\lambda\right]\\
&\qquad=
\Pr(\mathcal E_{\mathrm{ret}}\mid\iota,\boldsymbol\lambda)
\,\mathcal R_{P_\iota(\boldsymbol\lambda)}(T)
\le\mathcal R_{\max}(T).
\end{aligned}
\]
On failure, the output is zero and \(0\le\tau_\iota\le1\), so
\[
\mathbb E\left[
(\widetilde T-\tau_\iota)^2\mathbf1_{\mathcal E_{\mathrm{ret}}^c}
\mid\iota,\boldsymbol\lambda\right]
\le\Pr(\mathcal E_{\mathrm{ret}}^c\mid\iota,\boldsymbol\lambda)
\le\frac{\Delta_{\mathrm{tar}}^2}{128}.
\]
Averaging under either latent prior gives
\[
\mathcal A_\iota(\widetilde T)
\le\mathcal R_{\max}(T)+\frac{\Delta_{\mathrm{tar}}^2}{128}.
\]
Combining this with the raw testing lower bound yields
\[
\frac{\Delta_{\mathrm{tar}}^2}{64}
\le\max_\iota\mathcal A_\iota(\widetilde T)
\le\mathcal R_{\max}(T)+\frac{\Delta_{\mathrm{tar}}^2}{128},
\qquad
\mathcal R_{\max}(T)\ge\frac{\Delta_{\mathrm{tar}}^2}{128}.
\]
% lean: CausalSmith.Stat.AnnotationRarearmFrontier.affine_fixed_rule_gap_lower

\item
For \(x\ge0\), the two cases \(x\le1\) and \(x>1\) give
\[
(\min\{x,1\})^2
=
\begin{cases}
x^2,&x\le1,\\
1,&x>1
\end{cases}
=\min\{1,x^2\}.
\]
Therefore every original clipped Borel randomized rule satisfies
\[
\mathcal R_{\max}(T)
\ge\frac{\Delta_{\mathrm{tar}}^2}{128}
\ge\frac{10^{-20}}{128}\min\{1,x^2\}.
\]
The constant zero rule belongs to the rule class. Taking the infimum over that class, as specified in \cref{obj:def:risk}, proves
\[
R(n,m,d,\epsilon)
\ge\frac{10^{-20}}{128}\min\{1,x^2\}.
\]
Together with the support and marginal equalities already established, this proves the assertion with the stated numerical constant \(c\).
% lean: CausalSmith.Stat.AnnotationRarearmFrontier.affine_fixed_minimax_rate_lower
\end{enumerate}
\end{proof}
